% header
\documentclass[openany, 14px]{book}
\usepackage{graphicx} % Required for inserting images
\usepackage{hyperref}

\title{\textbf{book}}
\author{Hrishik Yendluri}
\date{}

% document
\begin{document}
\fontsize{12pt}{18pt} \selectfont

% title page
\maketitle

% acknowledgments

% table of contents

% preface
\newpage
\chapter*{\textit{Preface}}
This book was written to help people new to the field computer science learn more about the low level aspects of programming through the implementation of an operating system. Operating systems are closely intertwined with the hardware of a computer and therefore can be an effective method for learning all things low level. This book does not take a direct approach to the implementation of the OS. Rather, it identifies the key concepts and explains how one can go about programming them to allow for originality.
\\\\
This book is written on my experience developing a (toy) operating system and is heavily based upon the chronological implementation of features within that personal OS. The OS comes with some documentation to explain how it was programmed, which can be found on my website (\url{https://hrishiky.github.io/static/projects/hrishikos/index.html}). The operating system source code can be found in the following repository: \url{https://github.com/hrishiky/hrishikos}. The LaTeX code for this book can be found in the following repository: \url{INSERT LINK}. If this book or the reference operating system helped, please consider starring the repositories.
\clearpage

% prerequisites
\newpage
\chapter*{0 Prerequisites}
\setcounter{section}{-1}
\clearpage

\newpage
\section{Introduction}
This section will discuss and explain the prerequisites needed for writing an operating system, including both skills and knowledge. The first two sections after this one will discuss the soft skills and the sections afterward will discuss the computer science knowledge needed.

\section{Mindset}

\section{Research}
\clearpage

\end{document}
